\section{Ensaio sobre o Totem}
\subsection{Introdução}

O pensamento de Mauss apresenta uma potência \textit{sui generis} que parece ter escapado ao olhar do próprio autor.
Classifico-a assim para evidenciar a associação que encontro entre as generalizações indicadas por ele em seus diversos trabalhos e o caráter que assume apresentar a eficácia da magia.
Sua obra prima, o \enquote{Ensaio sobre a Dádiva} (ED), parte de investigações sobre o direito contratual e o sistema de prestações econômicas em sociedades primitivas para encontrar uma forma contratual geral que consiste num sistema de prestações totais.
Composto por três obrigações -- \textit{dar, receber e retribuir} -- e tendo por princípio motor o \textit{hau} -- o espírito das coisas misturado ao das pessoas --, tal sistema foi classificado pelo autor como um \textit{fato social total} por abarcar ao mesmo tempo todas as instituições de uma dada sociedade:
religiosas, jurídicas, morais, econômicas, além de fenômenos de ordem estética, linguística, morfológica e social.
\enquote{Acreditamos ter aqui encontrado uma das rochas humanas sobre as quais são construídas nossas sociedades} (p. 187), dis o autor, cujos trabalhos tendem a se orientar pela busca de mecanismos os mais genéricos e gerais possíveis em relação à organização social de seres humanos através da observação de \enquote{sistemas o mais \textit{heterogêneos} possíveis} (p. 50) -- como não só o ED, mas também seu \enquote{Esboço de uma teoria geral da magia} (ETM), são ótimos exemplares.
Neste, ao contrário do que faz no ED, Mauss parte de uma noção vaga -- a de magia -- para encontrar-lhe uma definição consistente, resultando numa composição entre agentes, representações e \enquote{atos tradicionais de uma eficácia \textit{sui generis}} (p. 55), que possuem, no entanto, um caráter oculto, ou até oposto, ao coletivo, embora seu princípio motor -- o \textit{mana} -- derive sua força de crenças coletivas.
De fato, um certo determinismo do social parece transpassar o pensamento de Mauss funcionando como um fundamento para suas ideias.
Mesmo à noção de \textit{hau} é atribuída uma eficácia medida socialmente, visto, entre outras coisas, realizar-se através de cálculos valorativos acerca das dádivas em circulação, como as práticas do \textit{kula} e do \textit{potlatch} ilustram bem.
Encontra-se assim um problema em suas generalizações, pois sociedades humanas não podem prescindir dos seres humanos que as compõem.
Se as explicações dadas demandam sempre algum aspecto tradicional, todas se perderão se a gênese da tradição -- ou da sociedade -- em si for posta em questão.
Uma citação esclarecedora de \enquote{As técnicas do corpo} (TC) diz o seguinte:
\enquote{todos cometemos, e cometi durante muitos anos, o erro fundamental de só considerar que há técnica quando há instrumento (...). Chamo técnica um ato tradicional eficaz (...). Não há técnica nem transmissão se não houver tradição} (p. 405).
Escapa-lhe neste trecho que a tradição é também um instrumento, assim como a contradição que se verificaria na análise da primeira técnica adquirida por um ser humano.
No mesmo texto, sua hipótese de \enquote{homem total} -- composta por três dimensões:
biológica, psicológica e social -- torna a contradição ainda mais evidente, pois o social não esteve presente desde sempre:
ele se constrói historicamente tanto na escala da espécie humana quanto na do indivíduo.
Há, portanto, uma tautologia fundamental no pensamento maussiano:
o social gerando, movendo e reproduzindo o social.
O que me desconcerta é o modo como se encontra no conjunto de suas generalizações a resolução de sua própria tautologia:
esta é a potência \textit{sui generis} que vejo em seu pensamento.

Meu objetivo neste ensaio é indicar como Mauss, ao buscar os fundamentos que movem as sociedades humanas, encontrou um conjunto de padrões que -- apesar da redundância com que foram apresentados pelo autos -- mapeiam os princípios fundamentais e motores da própria existência.
Por este motivo, proponho uma releitura de seu \enquote{fato social total} como um \textit{fato \textbf{ontológico} total}, ancorada por um dispositivo analítico recursivo centrado numa estrutura que chamei de \textit{\textbf{Totem}}.

Buscarei descrever este dispositivo através da demonstração de seus movimentos em diferentes trabalhos maussianos, assim como na investigação de Caroline Garlet de Oliveira (2016) sobre o processo de aprendizado da dança Malinké ao longo das aulas oferecidas pelo Grupo Abayomi na Universidade Federal de Santa Catarina.
Dada a centralidade do corpo e do aprendizado de uma técnica para o seu trabalho, suas observações remontam num primeiro momento àquelas de Mauss em TC;
contudo, uma análise centrada no movimento de Totens revela articulações importantes com as demais obras já citadas, e também com a \enquote{Morfologia Social} (MS) e \enquote{Uma categoria do espírito humano: a noção de pessoa, e de \enquote{Eu}} (NPE).
Totods estes trabalhos revelam facetas distintas dessa estrutura e de outras categorias que se desdobram a partir dela, sendo a da \enquote{Existência} -- uma das mais fundamentais -- uma transfiguração do sistema de prestações totais que, ao remover o pressuposto tautológico do social, revela a riqueza do pensamento mágico para se pensar um princípio geral de causalidade -- um \textit{fato causal total}.

\subsection{O Totem}

A origem do Totem remonta a um ensaio de Bruno Latour chamado \enquote{\textit{\textbf{Como falar do corpo? A dimensão normativa dos estudos sobre a ciência}}} (2008), em que o autor define o corpo como uma \enquote{\textit{interface que vai ficando mais descritível quando aprende a ser afetado por muitos mais elementos}} a fim de superar dicotomias como as de corpo e mente, cultura e natureza, e objetividade e subjetividade, resultando num corpo que é \enquote{\textit{aquilo que deixa uma trajetória dinâmica através da qual aprendemos a registrar e a ser sernsíveis àquilo de que é feito o mundo}}.
Este aprendizado rejeita o uso de afirmações, dando-se através de proposições que emergem conforme o corpo se sensibilizar, isto é, aprende a ser afetado, a contrastes que se fazem progressivamente notáveis conforme sua exposição sucessiva a cenários sutilmente diferenciados.

O processo é ilustrado com uma descrição do treinamento de profissionais da indústria de perfumes que, cmo o recurso de um kit de odores \enquote{constituído por uma série de fragrâncias puras nitidamente distsintas, dispostas de forma a poder passar-se do contraste mais abrupto ao mais suave}, os profissionais aprendem a ser afetados por diferenças olfativas cada vez mais sutis.
Através do treinamento, o corpo adquire um nariz que lhe permite habitar num mundo mais diferenciado, levando à conclusão de que \enquote{as partes do corpo (...) são adquiridas progressivamente ao mesmo tempo que as «contrapartidas do mundo» vão sendo registradas de nova forma. \textit{Adquirir um corpo é um empreendimento progressivo que produz simultaneamente um meio sensorial e um mundo sensível.}}

É desta proposição que deriva a estrutura do Totem, e que se divide em três dimensões -- Mancha, Marca e Máscara -- correspondentes, grosso modo, ao meio sensorial, ao mundo sensível, e a este empreendimento progressivo de aquisição corporal.
Articulando esta estrutura às diversas generalizações trazidas por Mauss, encontrei um movimento análogo ao do sistema de trocas do ED capaz tanto de descrevê-las quanto àquilo que Latour entende como processo de aquisição de um corpo.
É a este movimento que me refiro pela categoria da Existência, um circuito recursivo em três atos -- Exibição, Reconhecimento e Mancha -- correspondentes às três obrigações do circuito da dádiva.
Apesar de haver uma diferença sensível entre a compreensão do princípio motor das trocas e o da Existência, dei a este o nome de \textit{Hau} a fim de demarcar a familiaridade entre as teorias, a qual buscarei aqui demonstrar.
Para tal, apresentarei a seguir a categoria do \textit{Ser}, que será fundamental para, por fim, debater a agenda política de Mauss à luz do Totem.

\subsection{Marcas que Mancham o Ser em movimento}

Dois contextos etnográficos são fundamentais para a descrição do sistema de prestação totais por Mauss:
as sociedades polinésias e as do noroeste americano.
Estas trariam o caso mais emblemático, o \textit{potlatch}, um sistema de prestações totais de tipo agonístico, marcado por profunda rivalidade entre seus participantes -- uma competição por prestígio -- e pela \enquote{\textit{obrigação absoluta de retribuir a dádiva}} (p. 193) a fim de garantir a autoridade em disputa.
Ali, o sistema opera via trocas de todo tipo de bens, pessoas, festas, serviços, etc, podendo-se chegar a batalhas mortais ou, até mesmo, à destruição suntuária de riquezas.
As prestações totais são conferidas pelo fato de que \enquote{é claramente o clã inteiro que contrata por todos, poro tudo o que ele possui e por tudo o que ele faz, mediante seu chefe} (p. 190), o que traz um caráter importante da instituição:
as trocas se dão entre pessoas morais -- clãs, grupos familiares e outros tipos de coletividades humanas.